Al terminar la práctica nos quedaron algunas preguntas, y estas son las respuestas que pudimos darles.

\textbf{1. ¿Por qué medir el camino desde A' da lo mismo que medirlo desde A?}

Porque A' es el reflejo de A. Cualquier punto del espejo queda a la misma distancia de los dos, como si el espejo fuera una bisagra que dobla el papel por la mitad.

\textbf{2. ¿Qué pasaría si el espejo estuviera dentro del agua?}

El camino más corto seguiría siendo el mismo, porque el rayo nunca cambia de medio y su velocidad, aunque menor, es igual en todo el recorrido. El tiempo total sería mayor, pero el punto donde ese tiempo es mínimo no se movería.

\textbf{3. ¿La luz siempre toma el camino de menor tiempo?}

No exactamente. La versión más correcta del principio dice que el tiempo debe ser estacionario, o sea que su derivada valga cero, y eso incluye mínimos, pero también máximos o puntos de inflexión. En espejos curvos, por ejemplo, la luz puede seguir caminos de tiempo máximo.

\textbf{4. ¿Cómo se podría repetir el experimento con refracción?}

Habría que dibujar una línea que separe dos medios y, en lugar de medir solo longitudes, multiplicar cada tramo por el índice de refracción de su medio para obtener el camino óptico. El mínimo de esa cantidad debería cumplir la ley de Snell de la ecuación \eqref{eq:snell}.
